\documentclass[10pt,a4paper]{amsart}
\usepackage[margin=19mm]{geometry}
\usepackage{amsmath,amssymb,mathtools}
\usepackage[scr=rsfs,cal=euler]{mathalfa}
\usepackage{xfrac}
\usepackage[unicode,hidelinks,bookmarks=false]{hyperref}
\newcommand{\ie}{i.e.\ }
\newcommand{\cf}{cf.\ }
\newcommand{\resp}{respectively\ }
\newcommand{\Subsection}{\subsection}
\providecommand{\nobreakdash}{\nobreak\hbox{-}}
\setlength{\emergencystretch}{2em}
\multlinegap0pt
\usepackage{amssymb}
\usepackage{dsfont}
\usepackage{bm}
\usepackage{xcolor}
\usepackage{tikz}
\newtheorem{lemma}{Lemma}
\newcommand{\Cnt}{\mathbf{Count}}
\newcommand{\NextPos}{\mathbf{NextPos}}
\newcommand{\rem}{\mathrm{rem}}
\newcommand{\subseq}{\mathrm{subseq}}
\newcommand{\unrankk}{\operatorname{stringUnrank}}
\title{Improved Capacity Upper Bounds for the Deletion Channel using a Parallelized Blahut-Arimoto Algorithm}
\author{Martim Pinto, Jo\~ao Ribeiro}
\date{Source: arXiv:2604.05867v1 (CC BY 4.0)}
\begin{document}
\maketitle

\noindent\emph{Source and licence.} Improved Capacity Upper Bounds for the Deletion Channel using a Parallelized Blahut-Arimoto Algorithm. Version arXiv:2604.05867v1 (CC BY 4.0), distributed by arXiv under the Creative Commons Attribution 4.0 International licence, http://creativecommons.org/licenses/by/4.0/, as recorded in the arXiv licence metadata for this exact version and shown on the abstract page https://arxiv.org/abs/2604.05867v1. Full licence notice: \texttt{licenses/arxiv-2604.05867-CC-BY-4.0.txt}.\par\smallskip

\noindent\textit{Editorial scope.} Counted unit: the article's lemma lem:surj (source0.tex lines 672--679). The article's complete original proof of it is delivered with it (source0.tex lines 681--743, one \texttt{proof} environment of 63 lines). No mathematics is written, repaired or re-derived: the statement and the proof are the article's own lines, in the article's own notation, and no retained line is substituted or re-flowed.  No further source-local context is needed: every cross-reference of every retained line resolves inside the two delivered spans.  Nothing was omitted from any retained span, so the package carries no omission record.  No retained line contains a citation, so the wrapper carries no bibliography and no bibliography entry.  No retained line contains a figure environment, an image-inclusion command or drawing code, so no image is delivered and the package declares no asset dependency.  Editorial formatting only, in addition: an AMS \texttt{amsart} preamble that restates the article's own declarations and macros used by the retained lines, namely the environments \texttt{lemma} on separate counters, together with the standard packages those lines need.  The retained environments print in this document as Lemma 1 in order of appearance, whereas the article prints the same items with the numbering of its own full document.  The complete native source of the article as distributed in the arXiv e-print for this exact version is bundled unchanged as \texttt{sources/197983/source0.tex}.
\begin{lemma}\label{lem:surj}
If $0\le j<N_{x,k}$, then $\unrankk(x,k,j)$ returns a length-$k$ subsequence of $x$. 
More precisely, at the start of iteration $t$ (i.e., after $t$ bits have been appended), the invariant
\[
0 \le j < \Cnt[i][\rem],\qquad \text{where }\rem=k-t
\]
holds, and $\subseq$ is a subsequence of $x_{[0: i-1]}$.
\end{lemma}

\begin{proof}
We prove this statement by induction.

\paragraph{Initialization.} At the start of the execution of the algorithm we have $t=0$, $\rem=k$, $i=0$, and the precondition requires $0\leq j<N_{x,k}=\Cnt[0][k]$.  
The empty string $\subseq$ is trivially a subsequence of the empty prefix.


\paragraph{Inductive step.}  
Assume the invariant holds at the $t$-th iteration, i.e.,
\[
    0 \le j < \Cnt[i][\rem],
    \quad \subseq \subseteq x_{[0: i-1]},
    \quad \rem=k-t>0.
\]
Let
\[
j_0 = \NextPos[i][0], 
\qquad 
c_0 = \Cnt[j_0+1][\rem-1],
\]
with $c_0=0$ if $j_0=n$.  
Similarly, let
\[
j_1 = \NextPos[i][1], 
\qquad 
c_1 = \Cnt[j_1+1][\rem-1],
\]
with $c_1=0$ if $j_1=n$.  
By the recurrence relation for $\Cnt$ we have
\[
\Cnt[i][\rem] = c_0 + c_1.
\]

We proceed by cases.
\begin{enumerate}
    \item \emph{($j<c_0$)} In this case the algorithm appends $0$ at position $j_0$, sets $i\gets j_0+1$, $\rem\gets\rem-1$, and leaves $j$ unchanged.  
    Since $j<c_0=\Cnt[j_0+1][\rem-1]$, we obtain
    \[
    0 \le j < \Cnt[i][\rem],
    \]
    with the updated $i,\rem$. The new subsequence is valid because it extends by a $0$ occurring at $j_0$.

    \item \emph{($j\ge c_0$)}  
In this case we update $j\gets j-c_0$. From the recurrence
\[
\Cnt[i][\rem]=c_0+c_1,
\]
the condition $j<\Cnt[i][\rem]$ implies $j-c_0<c_1$.  
The algorithm appends $1$ at position $j_1$, sets $i\gets j_1+1$, $\rem\gets\rem-1$.  
Thus the invariant becomes
\[
0 \le j < \Cnt[i][\rem],
\]
with the new $i,\rem$, and the subsequence remains valid.
\end{enumerate}  


In both cases the invariant is preserved.

\medskip\noindent
\textbf{Termination.}  
Each loop decreases $\rem$ by one. After $k$ iterations we have $\rem=0$ and exactly $k$ bits appended. By the invariant, $\subseq$ is a subsequence of $x$ of length $k$. Thus the algorithm always returns a valid subsequence of length $k$.
\end{proof}

\end{document}
